\newglossaryentry{slicing}{
    name={Slicing},
    description={
            Slicing bezeichnet die algorithmische Zerlegung eines dreidimensionalen Modells in diskrete Schichten sowie die Ableitung der zugehörigen Werkzeugpfade für den Fertigungsprozess.
    }
}

\newglossaryentry{infill}{
    name={Infill},
    description={
            Infill beschreibt die innere Füllstruktur eines additiv gefertigten Bauteils, die mechanische Eigenschaften, Fertigungszeit und Materialverbrauch maßgeblich beeinflusst.
    }
}

\newglossaryentry{inverse-kinematics}{
    name={Inverse Kinematik},
    description={
            Die inverse Kinematik umfasst die Berechnung notwendiger Gelenkkoordinaten aus einer vorgegebenen kartesischen Zielpose des Werkzeugs und ist zentral für die Bahnplanung von Robotersystemen.
    }
}

\newglossaryentry{non-coplanar}{
    name={nicht-koplanar},
    description={
            Nicht-koplanare Fertigungsstrategien nutzen Schicht- oder Bahnverläufe, die nicht in einer gemeinsamen Ebene liegen, um geometrische Freiheitsgrade des Prozesses besser auszuschöpfen.
    }
}

\newglossaryentry{anisotropic-component-properties}{
    name={anisotrope Bauteileigenschaften},
    description={
            Anisotrope Bauteileigenschaften beschreiben richtungsabhängige mechanische, thermische oder strukturelle Kennwerte eines additiv gefertigten Bauteils infolge der schichtweisen Fertigung.
    }
}

\newglossaryentry{medusa}{
    name={MEDUSA},
    description={
            MEDUSA ist die im Rahmen dieser Arbeit prototypisch entwickelte Softwarekomponente für nichtplanares \gls{slicing} und Pfadplanung in der robotergestützten additiven Fertigung.
    }
}

\newglossaryentry{toolpath}{
    name={Toolpath},
    plural={Toolpaths},
    description={
            Ein Toolpath bezeichnet die chronologisch geordnete Folge von Bewegungssegmenten, die das Fertigungswerkzeug während des additiven Aufbauprozesses abfährt.
    }
}

\newglossaryentry{nonplanar-slicing}{
    name={nichtplanares Slicing},
    description={
            Nichtplanares Slicing bezeichnet Schichtzerlegungsverfahren, bei denen die einzelnen Fertigungsschichten nicht in parallelen Ebenen, sondern entlang gekrümmter oder gekippter Flächen verlaufen.
    }
}

\newglossaryentry{singularity}{
    name={Singularität},
    plural={Singularitäten},
    description={
            Eine Singularität bezeichnet eine Konfiguration eines Robotersystems, in der die kinematische Kette einen oder mehrere Freiheitsgrade verliert und die Jacobi-Matrix singulär wird.
    }
}

\newglossaryentry{postprocessor}{
    name={Postprozessor},
    plural={Postprozessoren},
    description={
            Ein Postprozessor transformiert maschinenunabhängige Werkzeugpfaddaten in steuerungsspezifische Bewegungsbefehle für eine konkrete Fertigungsanlage.
    }
}

\newglossaryentry{kronos}{
    name={KRONOS},
    description={
            KRONOS ist die im Rahmen dieser Arbeit konzipierte Softwarekomponente für die roboterkinematische Planung und Echtzeitansteuerung des 6-Achs-Robotersystems in der additiven Fertigung.
    }
}

\newglossaryentry{additive-manufacturing}{
    name={Additive Fertigung},
    description={
            Die additive Fertigung umfasst Verfahren, bei denen ein Bauteil durch das gezielte, schichtweise Hinzufügen von Material auf Basis eines digitalen Modells aufgebaut wird.
            Sie steht im Gegensatz zu subtraktiven Verfahren, bei denen Material entfernt wird.
    }
}

\newglossaryentry{material-extrusion}{
    name={Materialextrusion},
    description={
            Materialextrusion bezeichnet ein additives Fertigungsverfahren, bei dem ein thermoplastischer Werkstoff durch eine beheizte Düse gefördert und schichtweise auf einer Bauplattform abgelegt wird.
    }
}

\newglossaryentry{cartesian-kinematics}{
    name={Kartesische Kinematik},
    description={
            Eine kartesische Kinematik beschreibt ein Bewegungssystem mit drei zueinander orthogonalen, translatorischen Achsen.
            Sie ist das vorherrschende Bauprinzip konventioneller 3D-Drucker.
    }
}

\newglossaryentry{tool-orientation}{
    name={Werkzeugorientierung},
    description={
            Die Werkzeugorientierung beschreibt die räumliche Ausrichtung des Fertigungswerkzeugs relativ zum Bauteil.
            In mehrachsigen Systemen ist sie ein zusätzlicher Freiheitsgrad gegenüber rein translatorischen Kinematiken.
    }
}

\newglossaryentry{path-planning}{
    name={Bahnplanung},
    description={
            Bahnplanung bezeichnet die Berechnung von Bewegungspfaden des Werkzeugs unter Berücksichtigung von Geometrie, Kollisionsfreiheit, kinematischen Restriktionen und prozesstechnischen Anforderungen.
    }
}

\newglossaryentry{degrees-of-freedom}{
    name={Freiheitsgrade},
    description={
            Die Freiheitsgrade eines Bewegungssystems geben die Anzahl der unabhängigen, gleichzeitig steuerbaren Bewegungen an.
            Ein 6-Achs-Roboter besitzt sechs Freiheitsgrade, drei translatorisch und drei rotatorisch.
    }
}

\newglossaryentry{6-axis-slicer}{
    name={6-Achs-Slicer},
    description={
            Ein 6-Achs-Slicer ist ein Slicing-System, das neben der Schichtzerlegung auch die Werkzeugorientierung entlang lokaler Geometrieeigenschaften plant und damit nicht-planare Schicht- und Bahnstrategien für 6-Achs-Robotersysteme ermöglicht.
    }
}

\newglossaryentry{mesh}{
    name={Mesh},
    plural={Meshes},
    description={
            Polygonale Oberflächendarstellung eines dreidimensionalen Körpers, in dieser Arbeit ausschließlich als Dreiecksnetz aus \glspl{vertex}, \glspl{edge} und \glspl{triangle}.
            Ein Mesh dient als Eingabe der Slicing-Algorithmen.
    }
}

\newglossaryentry{vertex}{
    name={Vertex},
    plural={Vertizes},
    description={
            Eckpunkt eines \gls{mesh}.
            Jeder Vertex ist durch seine Position im dreidimensionalen Raum sowie optional durch eine Normale gekennzeichnet, die für die Werkzeug-Anstellrichtung in nichtplanaren Slicing-Verfahren interpoliert wird.
    }
}

\newglossaryentry{edge}{
    name={Kante},
    plural={Kanten},
    description={
            Verbindung zwischen zwei \glspl{vertex} eines \gls{mesh}.
            In einem \gls{watertight}[en] \gls{mesh} gehört jede Kante genau zu zwei \glspl{triangle}.
            Kanten sind die elementaren Träger der Schnittinformation: ein Schnitt entsteht dort, wo die Endpunkte einer Kante auf verschiedenen Seiten der \gls{slice-plane} liegen.
    }
}

\newglossaryentry{triangle}{
    name={Dreieck},
    plural={Dreiecke},
    description={
            Polygonale Grundfläche eines \gls{mesh}, definiert durch drei \glspl{vertex} und drei \glspl{edge}.
            Bei der Schnittoperation liefert ein Dreieck unter regulären Voraussetzungen höchstens ein \gls{line-segment} pro \gls{slice-plane}.
    }
}

\newglossaryentry{watertight}{
    name={wasserdicht},
    description={
            Eigenschaft eines \gls{mesh}, vollständig geschlossen und mannigfaltig zu sein.
            Jede \gls{edge} gehört genau zu zwei \glspl{triangle}, sodass sich Schnittsegmente eindeutig zu geschlossenen Konturen verketten lassen.
            Wasserdichtheit ist Voraussetzung für die Korrektheit aller in dieser Arbeit umgesetzten Slicing-Algorithmen.
    }
}

\newglossaryentry{manifold}{
    name={Mannigfaltigkeit},
    plural={Mannigfaltigkeiten},
    description={
            Topologische Eigenschaft eines \gls{mesh}, in der jeder Punkt eine Umgebung besitzt, die homöomorph zu einer offenen Scheibe der Ebene ist.
            Jede \gls{edge} gehört zu höchstens zwei \glspl{triangle}.
            T-Stücke, Selbstdurchdringungen und doppelt belegte Kanten sind ausgeschlossen.
    }
}

\newglossaryentry{support-structures}{
    name={Stützstruktur},
    plural={Stützstrukturen},
    description={
            In der konventionellen \ac{FDM}-Fertigung temporär mitgedrucktes Hilfsmaterial, das überhängende Bauteilbereiche jenseits eines kritischen Überhangwinkels während der Fertigung abstützt und anschließend mechanisch oder chemisch entfernt wird.
    }
}

\newglossaryentry{build-axis}{
    name={Aufbauachse},
    plural={Aufbauachsen},
    description={
            Globale Richtung, entlang der ein Bauteil in der konventionellen additiven Fertigung schichtweise aufgebaut wird.
            Sie definiert in planaren Slicing-Verfahren die Normale aller \glspl{slice-plane} sowie die konstante \gls{tool-orientation}.
    }
}

\newglossaryentry{layer-height}{
    name={Schichtdicke},
    plural={Schichtdicken},
    description={
            Konstanter Höhenabstand zweier aufeinanderfolgender \glspl{slice-plane} entlang der \gls{build-axis} bzw.\ einer lokalen \gls{growth-direction}.
            Bestimmt die vertikale Auflösung des gefertigten Bauteils und stellt einen zentralen Parameter aller Slicing-Algorithmen dar.
    }
}

\newglossaryentry{growth-direction}{
    name={Wachstumsrichtung},
    plural={Wachstumsrichtungen},
    description={
            Lokale Richtung, entlang der ein einzelner Slicing-Lauf eines nichtplanaren Verfahrens Schichten aufbaut.
            Im Gegensatz zur global festgelegten \gls{build-axis} kann jeder \gls{run} eine eigene Wachstumsrichtung besitzen, was die Ausnutzung der zusätzlichen Freiheitsgrade einer 6-\ac{DOF}-Kinematik ermöglicht.
    }
}

\newglossaryentry{overhang}{
    name={Überhang},
    plural={Überhänge},
    description={
            Bauteilregion, deren Oberfläche stärker gegen die aktuelle \gls{growth-direction} geneigt ist als ein parametrierter Grenzwinkel und die deshalb in der konventionellen Fertigung \glspl{support-structures} erfordert.
            Bei nichtplanaren Slicing-Verfahren kann die Erkennung eines Überhangs einen Wechsel der \gls{growth-direction} auslösen und so eine stützfreie Fertigung ermöglichen.
    }
}

\newglossaryentry{run}{
    name={Lauf},
    plural={Läufe},
    description={
            Verarbeitungseinheit nichtplanarer Slicing-Verfahren mit dynamisch wechselnden Wachstumsrichtungen.
            Ein Lauf umfasst die Erzeugung aller Schichten eines zusammenhängenden Bauteilabschnitts entlang einer einzelnen \gls{growth-direction} und ist durch Wachstumsrichtung, Startposition sowie Zweigbezeichner eindeutig charakterisiert.
    }
}

\newglossaryentry{branch}{
    name={Zweig},
    plural={Zweige},
    description={
            Bei nichtplanaren Slicing-Verfahren mit dynamischem Richtungswechsel die Gesamtheit aller Schichten, die im selben \gls{run} erzeugt wurden.
            Zweige bilden eine baumartige Zerlegung des Bauteils, deren Wurzel der von der \gls{buildplate} ausgehende \gls{trunk} ist.
    }
}

\newglossaryentry{trunk}{
    name={Stamm},
    plural={Stämme},
    description={
            Initialer \gls{branch} eines nichtplanaren Slicing-Verfahrens, der von der \gls{buildplate} entlang der globalen \gls{build-axis} aufgebaut wird und aus dem alle weiteren \glspl{branch} hervorgehen.
    }
}

\newglossaryentry{contour-ring}{
    name={Konturring},
    plural={Konturringe},
    description={
            Geschlossener \gls{polygonal-chain}, der die Außenkontur eines Bauteilquerschnitts auf einer Schicht beschreibt.
            Konturringe entstehen typischerweise durch Verkettung der Schnittsegmente einer \gls{slice-plane} und werden in der Pfadplanung als zusammenhängender Bewegungsbefehl ausgeführt.
    }
}

\newglossaryentry{liang-barsky}{
    name={Liang-Barsky-Verfahren},
    description={
            Parametrisches Verfahren zum Clipping von Liniensegmenten an einem achsenparallelen Rechteck.
            Das Verfahren berechnet die Eintritts- und Austrittsparameter entlang des Segments und liefert das beschnittene Segment ohne iterative Außen-Innen-Tests.
    }
}

\newglossaryentry{gram-schmidt}{
    name={Gram-Schmidt-Verfahren},
    description={
            Verfahren zur Konstruktion eines orthonormalen Basistripels aus einer linear unabhängigen Ausgangsmenge.
            In nichtplanaren Slicing-Verfahren erzeugt eine Gram-Schmidt-artige Konstruktion aus der \gls{growth-direction} und einer Referenzrichtung das lokale Koordinatensystem $(\vec{U}, \vec{V}, \vec{G})$.
    }
}

\newglossaryentry{snap-tolerance}{
    name={Snap-Toleranz},
    plural={Snap-Toleranzen},
    description={
            Numerische Toleranzumgebung im Bereich weniger Mikrometer, mit der bei der Verkettung von Schnittsegmenten geprüft wird, ob zwei berechnete Punkte als identisch zu betrachten sind.
            Sie überdeckt Rundungsfehler der \gls{ieee-754}-Arithmetik, ohne tatsächlich verschiedene Punkte zu verschmelzen.
    }
}

\newglossaryentry{bounding-rectangle}{
    name={Begrenzungsrechteck},
    plural={Begrenzungsrechtecke},
    description={
            Kleinstes achsenparalleles Rechteck, das eine zweidimensionale Punktmenge vollständig einschließt.
            Bei der Überhangerkennung in nichtplanaren Slicing-Verfahren wird pro Schicht das Begrenzungsrechteck in der \gls{uv-plane} berechnet und mit dem der Vorgängerschicht verglichen.
    }
}

\newglossaryentry{uv-plane}{
    name={UV-Ebene},
    plural={UV-Ebenen},
    description={
            Zweidimensionale \gls{slice-plane} eines Slicing-Laufs, aufgespannt durch die orthonormalen Achsen $\vec{U}$ und $\vec{V}$ senkrecht zur \gls{growth-direction}.
            In ihr werden die geometrische Überhangerkennung und das Beschneiden der Segmente durchgeführt.
    }
}

\newglossaryentry{buildplate}{
    name={Buildplate},
    plural={Buildplates},
    description={
            Bauplattform, auf der das Bauteil schichtweise aufgebaut wird.
            Sie bildet den Ursprung des initialen \gls{run} entlang der globalen \gls{build-axis}.
    }
}

\newglossaryentry{polygonal-chain}{
    name={Polygonzug},
    plural={Polygonzüge},
    description={
            Geordnete Folge von \glspl{line-segment}, deren Endpunkte jeweils mit dem Anfangspunkt des nachfolgenden Segments übereinstimmen.
            Ein geschlossener Polygonzug, in dem zusätzlich der letzte End- mit dem ersten Anfangspunkt zusammenfällt, bildet einen \gls{contour-ring}.
    }
}

\newglossaryentry{coverage}{
    name={Coverage},
    description={
            Vollständigkeitsmaß für nichtplanare Slicing-Verfahren, das den Anteil der durch Schichten abgedeckten Bauteilhöhe beschreibt.
            Werte nahe eins deuten auf eine vollständige Überdeckung hin, kleinere Werte auf unbearbeitete Regionen, etwa infolge erreichter Iterationslimits.
    }
}

\newglossaryentry{slice-plane}{
    name={Schnitt\-ebene},
    plural={Schnitt\-ebenen},
    description={
            Geometrische Ebene, mit der ein \gls{mesh} bei einer Schicht\-erzeugung verschnitten wird.
            Ihre Normale entspricht der \gls{build-axis} bei planaren bzw.\ der \gls{growth-direction} bei nichtplanaren Slicing-Verfahren; ihre Lage definiert die Position der erzeugten Schicht.
    }
}

\newglossaryentry{line-segment}{
    name={Liniensegment},
    plural={Liniensegmente},
    description={
            Durch zwei Endpunkte eindeutig festgelegtes, gerades Teilstück einer Geraden.
            Liniensegmente entstehen als Ergebnis der Schnittoperation eines \gls{triangle} mit einer \gls{slice-plane} und werden anschließend zu einem \gls{polygonal-chain} oder \gls{contour-ring} verkettet.
    }
}

\newglossaryentry{traveling-salesman-problem}{
    name={Traveling Salesman Problem},
    description={
            Klassisches Optimierungsproblem der kombinatorischen Optimierung, bei dem die kürzeste geschlossene Rundreise durch eine endliche Menge von Knoten gesucht wird.
            Im Slicing dient es als Referenzmodell für die Reihenfolgenoptimierung der \glspl{contour-ring} innerhalb einer Schicht.
            Aufgrund seiner NP-Schwere wird in der Praxis meist auf heuristische Verfahren wie das \gls{nearest-neighbor} zurückgegriffen.
    }
}

\newglossaryentry{nearest-neighbor}{
    name={Nächste-Nachbar-Verfahren},
    description={
            Gieriges Heuristikverfahren, das ausgehend von einem Startelement stets das in einem definierten Abstandsmaß nächstgelegene, noch nicht besuchte Element auswählt.
            Es wird in dieser Arbeit sowohl bei der Verkettung von \glspl{line-segment} zu \glspl{contour-ring} als auch bei der Reihenfolgenoptimierung der Konturringe innerhalb einer Schicht eingesetzt.
            Das Verfahren liefert keine Optimalitätsgarantie, ist aber für die hier auftretenden Problemgrößen praktisch ausreichend.
    }
}

\newglossaryentry{ieee-754}{
    name={IEEE 754},
    description={
            Internationaler Standard zur binären Darstellung von Gleitkommazahlen in Rechnern, der unter anderem die Datentypen \texttt{float} und \texttt{double} der Programmiersprache C++ definiert.
            Aufgrund der endlichen Mantissenlänge entstehen Rundungsfehler, die bei der Verkettung von Schnittsegmenten durch eine \gls{snap-tolerance} ausgeglichen werden müssen.
    }
}

\newglossaryentry{iso-surface}{
    name={Isofläche},
    plural={Isoflächen},
    description={
            Niveaufläche eines Skalarfeldes \(\Phi: \mathcal{M} \to \mathbb{R}\), entlang derer \(\Phi\) einen konstanten Wert \(\varphi^*\) annimmt, formal \(\{\vec{x} \in \mathcal{M} \mid \Phi(\vec{x}) = \varphi^*\}\).
            Isoflächen verallgemeinern den Begriff der Höhenlinie auf höhere Dimensionen und stehen lokal senkrecht auf dem Gradienten \(\nabla\Phi\).
            Ist \(\Phi\) hinreichend glatt und \(\varphi^*\) ein regulärer Wert (d.\,h. \(\nabla\Phi \neq 0\) auf der Niveaumenge), so bildet die Isofläche eine differenzierbare Untermannigfaltigkeit der Kodimension eins.
    }
}

\newglossaryentry{slice-surface}{
    name={Schnittfläche},
    plural={Schnittflächen},
    description={
            Verallgemeinerung der \gls{slice-plane} auf beliebig gekrümmte Flächen.
            Eine Schnittfläche zerlegt das Eingabe-\gls{mesh} in eine geordnete Folge von Schichten, ohne notwendigerweise eben zu sein.
            Im planaren und im Base-Algorithmus sind Schnittflächen stets eben, im Kristall-Algorithmus sind es gekrümmte \glspl{iso-surface}.
    }
}

\newglossaryentry{iso-step}{
    name={Iso-Schritt},
    plural={Iso-Schritte},
    description={
            Ein einzelner Schritt der adaptiven Iso-Flächen-Extraktion im Kristall-Algorithmus.
            Pro Iso-Schritt wird ein Iso-Wert \(\varphi^*_i\) gewählt, das \gls{mesh} mit der zugehörigen \gls{iso-surface} geschnitten und die entstehenden \glspl{line-segment} pro zusammenhängender Komponente zu einer Schicht verkettet.
            Der Index des Iso-Schritts wird in der Schichtrepräsentation mitgeführt und ist Grundlage des nachgelagerten Branch-Tracking sowie des Branch-Scheduling.
    }
}

\newglossaryentry{cotangent-laplacian}{
    name={Kotangens-Laplace-Operator},
    plural={Kotangens-Laplace-Operatoren},
    description={
            Diskretes Analogon des kontinuierlichen Laplace-Beltrami-Operators auf einem Dreiecksgitter.
            Die Außerdiagonalelemente sind durch \(L_{ij} = -\tfrac{1}{2}(\cot\alpha_{ij} + \cot\beta_{ij})\) gegeben, wobei \(\alpha_{ij}\) und \(\beta_{ij}\) die der \gls{edge} \((i,j)\) gegenüberliegenden Innenwinkel der beiden angrenzenden \glspl{triangle} sind.
            Der Operator ist symmetrisch, dünnbesetzt und konvergiert für quasi-uniforme Gitter gegen den kontinuierlichen Laplace-Beltrami-Operator.
    }
}

\newglossaryentry{quaternion}{
    name={Quaternion},
    plural={Quaternionen},
    description={
            Element des vierdimensionalen Schiefkörpers \(\mathbb{H}\), das eine Rotation im dreidimensionalen Raum singularitätsfrei beschreibt.
            Eine Einheits-Quaternion \((w, x, y, z)\) mit \(w^2 + x^2 + y^2 + z^2 = 1\) repräsentiert eine Drehung um eine Achse \((x, y, z)/\sqrt{1-w^2}\) mit Winkel \(2\arccos(w)\).
            Quaternionen vermeiden den Gimbal-Lock-Effekt von Eulerwinkeln und sind die Standardrepräsentation der Werkzeugorientierung in ROS und MoveIt.
    }
}

\newglossaryentry{idw}{
    name={Inverse Distance Weighting},
    description={
            Verfahren zur räumlichen Interpolation eines Wertes an einer Zielposition aus einer Menge bekannter Stützstellen.
            Der Wert ergibt sich als gewichtetes Mittel der Stützstellenwerte, wobei das Gewicht jeder Stützstelle invers proportional zu einer Potenz ihres Abstands zur Zielposition ist.
            Im Kristall-Algorithmus wird Inverse Distance Weighting verwendet, um Infill-Punkte aus der Tangentialebene auf die gekrümmte \gls{iso-surface} anzuheben.
    }
}

\newglossaryentry{pca}{
    name={Hauptkomponentenanalyse},
    plural={Hauptkomponentenanalysen},
    description={
            Statistisches Verfahren zur Bestimmung der Hauptachsen einer Punktwolke durch Eigenwertzerlegung der Kovarianzmatrix.
            Die Eigenvektoren mit den größten Eigenwerten geben die Richtungen maximaler Varianz an.
            Im Kristall-Algorithmus dient die Hauptkomponentenanalyse dazu, die Tangentialbasis einer Iso-Schicht so auszurichten, dass die Scanlinien des nicht-planaren Infill entlang der ausgedehnten Komponentenrichtung verlaufen.
    }
}

\newglossaryentry{infill-spacing}{
    name={Bahnabstand},
    plural={Bahnabstände},
    description={
            Senkrechter Abstand zwischen zwei benachbarten Scanlinien des \gls{infill}.
            Der Bahnabstand bestimmt die Dichte des Infill und damit das Verhältnis aus Materialverbrauch, Druckdauer und mechanischer Steifigkeit des gefertigten Bauteils.
    }
}

\newglossaryentry{even-odd-rule}{
    name={Even-Odd-Regel},
    plural={Even-Odd-Regeln},
    description={
            Punkt-in-Polygon-Test, der die Innen- und Außenbereiche eines Polygons über die Anzahl der Schnittpunkte eines beliebigen Strahls mit dem Polygonrand bestimmt.
            Eine ungerade Schnittzahl klassifiziert den Ausgangspunkt als innen, eine gerade als außen.
            Im Kristall-Algorithmus wird die Even-Odd-Regel verwendet, um die Schnittpunkte einer Scanlinie mit dem projizierten Konturpolygon paarweise zu Innen-Linienstücken zu gruppieren.
    }
}

\newglossaryentry{snake-pattern}{
    name={Mäanderbahn},
    plural={Mäanderbahnen},
    description={
            Anordnung paralleler Bahnen, bei der die Bewegungsrichtung jeder zweiten Bahn umgekehrt wird, sodass das Ende einer Bahn unmittelbar in den Anfang der nächsten übergeht.
            Im \gls{infill} vermeidet die Mäanderbahn Verfahrwege zwischen aufeinanderfolgenden Scanlinien und reduziert dadurch die Druckdauer sowie die Anzahl materialfreier Bewegungen.
    }
}

\newglossaryentry{heat-method}{
    name={Heat Method},
    description={
            Numerisches Verfahren zur Approximation geodätischer Distanzen auf einer Mannigfaltigkeit.
            Eine Wärmediffusionsgleichung wird über einen kurzen Zeitschritt gelöst, der normalisierte Wärmegradient liefert die Näherung der geodätischen Distanzrichtung.
            Die Heat Method wurde im Kristall-Algorithmus als Alternative zur harmonischen Feldformulierung erwogen, jedoch wegen ihrer Abhängigkeit von einem ad-hoc gewählten Diffusionszeitparameter verworfen.
    }
}

\newglossaryentry{drift-threshold}{
    name={Drift-Schranke},
    plural={Drift-Schranken},
    description={
            Maximaler räumlicher Abstand, den der Schwerpunkt einer Iso-Komponente zwischen zwei aufeinanderfolgenden \glspl{iso-step} aufweisen darf, um noch derselben \gls{branch} zugeordnet zu werden.
            Größere Sprünge werden im Branch-Tracking als Topologieereignis (Birth, Death, Split oder Merge) interpretiert.
            Die Drift-Schranke wird typischerweise als Vielfaches der \gls{layer-height} mit einer absoluten Untergrenze parametriert.
    }
}

\newglossaryentry{travel-move}{
    name={Verfahrweg},
    plural={Verfahrwege},
    description={
            Nicht-extrudierende Bewegung des Werkzeugs zwischen zwei Druckbewegungen.
            Verfahrwege treten innerhalb einer Schicht zwischen disjunkten Konturen, beim Wechsel zwischen \glspl{branch} und beim Übergang zwischen aufeinanderfolgenden Schichten auf.
            Sie werden in der \gls{toolpath}-Konstruktion explizit markiert, damit der Postprozessor in \ac{kronos} sie mit anderen Geschwindigkeitsprofilen und Extruderparametern versieht als die Druckbewegungen.
    }
}

\newglossaryentry{strategy-pattern}{
    name={Strategie-Entwurfsmuster},
    plural={Strategie-Entwurfsmuster},
    description={
            Verhaltens-Entwurfsmuster der objektorientierten Programmierung, das einen Algorithmus hinter einer abstrakten Schnittstelle kapselt und zur Laufzeit gegen alternative Implementierungen austauschbar macht.
            Im Kristall-Algorithmus wird das Branch-Scheduling als Strategie-Entwurfsmuster modelliert, sodass alternative Reihenfolgestrategien ohne Änderung der übrigen Pipeline integrierbar sind.
    }
}

\newglossaryentry{waypoint}{
    name={Wegpunkt},
    plural={Wegpunkte},
    description={
            Einzelner Punkt einer Trajektorie, der die vollständige 6-\ac{DOF}-Pose des Werkzeugs (Position und Orientierung) sowie die zugehörigen Prozessparameter (Vorschub, Extrusionsrate, Bewegungsart) zu einem definierten Zeitpunkt der Bewegungsfolge festlegt.
            Eine geordnete Folge von Wegpunkten definiert implizit die Segmente der Bewegungsfolge.
    }
}

\newglossaryentry{semver}{
    name={Semantic Versioning},
    description={
            Versionierungsschema in der Form \emph{MAJOR.MINOR.PATCH}, bei dem die Hauptversion (\emph{MAJOR}) bei nicht-abwärtskompatiblen Änderungen, die Nebenversion (\emph{MINOR}) bei abwärts\-kompatiblen Erweiterungen und die Patchversion (\emph{PATCH}) bei reinen Fehlerkorrekturen erhöht wird.
            Erlaubt automatisierte Kompatibilitätsprüfungen zwischen unabhängig versionierten Komponenten.
    }
}

\newglossaryentry{ipc}{
    name={Interprozesskommunikation},
    plural={Interprozesskommunikationen},
    description={
            Sammelbegriff für Mechanismen, die den Datenaustausch zwischen mehreren gleichzeitig laufenden Betriebssystemprozessen ermöglichen.
            Beispiele sind Pipes, Sockets, Shared Memory und nachrichtenbasierte Middleware wie ROS-2-Topics.
    }
}

\newglossaryentry{json}{
    name={JSON},
    description={
            \ac{JSON} ist ein textbasiertes, sprachunabhängiges Datenaustauschformat zur strukturierten Repräsentation von Schlüssel-Wert-Paaren, Arrays und verschachtelten Objekten.
            Aufgrund seiner Lesbarkeit, der einfachen Parsbarkeit und der breiten Unterstützung in nahezu allen Programmiersprachen hat sich \ac{JSON} als De-facto-Standard für den Datenaustausch zwischen heterogenen Softwarekomponenten etabliert.
            Im Rahmen dieser Arbeit dient \ac{JSON} als Übertragungsformat des geslicten Pfades zwischen \ac{medusa} und \ac{kronos}.
    }
}
\newglossaryentry{json-schema}{
    name={JSON-Schema},
    plural={JSON-Schemata},
    description={
            Formale Spezifikationssprache zur Beschreibung der Struktur und der zulässigen Werte von \gls{json}-Dokumenten.
            Erlaubt die maschinelle Validierung eingehender Dokumente gegen einen vorgegebenen Vertrag und damit die explizite Versionierung von Schnittstellen.
    }
}

\newglossaryentry{sphere-shell}{
    name={Sphärenschale},
    plural={Sphärenschalen},
    description={
        Zweidimensionale Mannigfaltigkeit aller Punkte des dreidimensionalen Raums, die von einem festen Mittelpunkt denselben Abstand $r$ besitzen.
        Geometrisch beschreibt sie die Oberfläche einer Kugel mit Radius $r$.
    }
}

\newglossaryentry{sphere-center}{
    name={Sphärenmittelpunkt},
    plural={Sphärenmittelpunkte},
    description={
        Punkt im dreidimensionalen Raum, von dem aus eine Sphäre oder eine Schar konzentrischer Sphärenschalen durch einen einheitlichen Radius definiert wird.
    }
}

\newglossaryentry{sphere-radius}{
    name={Sphärenradius},
    plural={Sphärenradien},
    description={
        Euklidischer Abstand zwischen dem Sphärenmittelpunkt und einem beliebigen Punkt auf der zugehörigen Sphärenschale.
        Charakterisiert zusammen mit dem Sphärenmittelpunkt eine Sphäre eindeutig.
    }
}

\newglossaryentry{ray-casting}{
    name={Ray-Casting},
    plural={Ray-Casting-Verfahren},
    description={
        Geometrisches Verfahren, bei dem von einem Testpunkt aus ein Strahl in eine beliebige Richtung ausgesandt und mit den Flächen einer Geometrie verschnitten wird.
        Die Anzahl bzw. Anordnung der Schnitte erlaubt Aussagen über Sichtbarkeit, Zugehörigkeit zu Volumina oder die Lage des Testpunkts relativ zur Geometrie.
    }
}

\newglossaryentry{ci-pipeline}{
    name={CI-Pipeline},
    plural={CI-Pipelines},
    description={
            Automatisierte Abfolge von Build-, Test- und Auslieferungsschritten, die bei jeder Änderung am Quellcode-Repository ausgeführt wird und die Grundlage der \gls{ci} bildet.
        }
}

\newglossaryentry{ci}{
    name={Continuous Integration},
    description={
            Praxis der Softwareentwicklung, bei der Quellcodeänderungen automatisch in einen gemeinsamen Hauptzweig integriert werden und dort durch automatisierte Build- und Testläufe kontinuierlich verifiziert werden.
            Ziel ist es, Integrationsprobleme frühzeitig zu erkennen und einen jederzeit lauffähigen Stand des Systems sicherzustellen.
        }
}

\newglossaryentry{family}{
    name={Schar},
    plural={Scharen},
    description={
        In der Geometrie eine durch einen oder mehrere Parameter indizierte Menge gleichartig konstruierter geometrischer Objekte wie Geraden, Ebenen, Kurven oder Flächen, deren Elemente sich ausschließlich in den jeweiligen Parameterwerten unterscheiden.
    }
}

\newglossaryentry{cmake}{
    name={CMake},
    description={
            CMake ist ein plattformübergreifendes Build-Konfigurationssystem, das aus deklarativen Beschreibungen native Build-Dateien für unterschiedliche Toolchains erzeugt und die Verwaltung von Targets, Sichtbarkeiten und Abhängigkeiten zentralisiert.
    }
}

\newglossaryentry{fetchcontent}{
    name={FetchContent},
    description={
            CMake-Modul zum Beziehen externer Quellprojekte zur Konfigurationszeit.
            Quelle und Versions-Tag werden deklarativ angegeben, anschließend wird das Projekt lokal ausgecheckt und als Teil des aktuellen Builds eingebunden, ohne dass eine systemweite Installation erforderlich ist.
    }
}

\newglossaryentry{FNV-1a}{
    name={FNV-1a-Hash},
    description={
            Nicht-kryptografischer Hash-Algorithmus aus der Familie der Fowler-Noll-Vo-Funktionen.
            Liefert einen 32- oder 64-Bit-Hashwert über einen Bytestrom durch alternierendes XOR und Multiplikation mit einer \ac{FNV}-Primzahl.
            Geeignet für Integritätsprüfungen und Hashtabellen, nicht jedoch für sicherheitskritische Anwendungen.
        }
}

\newglossaryentry{viridis}{
    name={Viridis},
    description={
            Wahrnehmungsuniforme Farbpalette mit monoton wachsender Helligkeit und ohne Bandbildung in den dominanten Farbsehschwächen.
            Sie wurde für die visuelle Darstellung kontinuierlicher Skalarfelder entwickelt und ist Standardpalette in zahlreichen wissenschaftlichen Visualisierungswerkzeugen.
    }
}

\newglossaryentry{face}{
    name={Face},
    plural={Faces},
    description={
        Elementare polygonale Zelle eines \gls{mesh}, die durch eine geordnete Folge von \glspl{vertex} aufgespannt wird und eine planare Teilfläche der diskreten Geometrie definiert.
        In Dreiecksgittern, die in der additiven Fertigung und im geometrischen Rechnen üblich sind, ist jeder Face ein Dreieck mit genau drei Eckpunkten.
        Benachbarte Faces teilen sich mindestens eine gemeinsame \gls{edge}.
        Jedem Face kann eine Flächennormale zugeordnet werden, die senkrecht auf der aufgespannten Ebene steht und die lokale Orientierung der Oberfläche kodiert.
    }
}

\newglossaryentry{segment}{
    name={Segment},
    plural={Segmente},
    description={
        Geradlinige Verbindung zwischen zwei Punkten im dreidimensionalen Raum, die eine elementare geometrische oder kinematische Einheit darstellt.
        Im geometrischen Kontext bezeichnet ein Segment das Ergebnis des Schnitts einer Ebene oder Iso-Fläche mit einem \gls{face}: ein gerichtetes Linienstück, dessen Endpunkte auf den \glspl{edge} des \glspl{face} liegen.
        Im Kontext der Werkzeugpfadplanung ist ein Segment die kleinste ausführbare Bewegungseinheit des Roboters; es trägt neben Anfangs- und Endpunkt eine Werkzeugorientierung sowie Metadaten wie Branch-Bezeichner, Schicht-Index und die Unterscheidung zwischen Druck- und Verbindungsbewegung.
        Aufeinanderfolgende Segmente bilden zusammen einen \gls{toolpath}.
    }
}

\makeglossaries